\documentclass[a5paper,10pt]{article}

% https://github.com/InDevRus/filippov-solutions

% Нумерация страниц
\pagenumbering{gobble}

% Настройка полей
\usepackage{geometry}
\geometry{left=1cm}
\geometry{right=1cm}
\geometry{top=1cm}
\geometry{bottom=1.5cm}

% Вставка таблицы
\usepackage{array}
\usepackage{tabularx}

% Инструменты выравнивания формул
\usepackage{amsmath}
\usepackage{mathtools}

% Диакритика в математике
\usepackage{amsfonts}

% Вычисления размеров на ходу
\usepackage{calc}

% Удобные записи производных
\usepackage[italic=true]{derivative}

% Настройки сносок
\usepackage{enotez}

% Длинные знаки векторов
\usepackage{anyfontsize}
\usepackage[b]{esvect}

% Вставка картинок
\usepackage{graphicx}

% Вставка красной строки
\usepackage{indentfirst}
\setlength{\parindent}{1em}

% Условия в командах
\usepackage{xifthen}

% Луа-скрипты
\usepackage{luacode}

% Настройка команд отдельно в математическом и текстовом режимах
\usepackage{mathcommand}

% Для повторения LaTeX кода
\usepackage{multido}

% Конфигурация отступов формул
\delimitershortfall-1sp
\usepackage{mleftright}
\mleftright

% Растягивание объектов
\usepackage{relsize}

% Обтекание текстом
\usepackage{wrapfig}
\usepackage{subcaption}
\usepackage{floatflt}

% Поддержка ввода кириллицы
\usepackage[english, russian]{babel}

% Настройка корневой позиции для компилятора
\usepackage{import}
\providecommand*{\ProjectRootPath}{..}

\import{\ProjectRootPath/preambles/macros/}{theorems}

\import{\ProjectRootPath/preambles/fonts}{font_setup}

% Знак градуса
\usepackage{siunitx}

\import{\ProjectRootPath/preambles/macros/}{operators}
\import{\ProjectRootPath/preambles/macros/}{redefinitions}

\import{\ProjectRootPath/preambles/fonts}{letters_setup}
\import{\ProjectRootPath/preambles/fonts/adjustments}{custom_superscripts}

\import{\ProjectRootPath/preambles/custom}{formatting}
\import{\ProjectRootPath/preambles/custom}{frames}
\import{\ProjectRootPath/preambles/custom}{list_setup}

% TODO: перевести команды на современный стиль объявления

\begin{document}
    $\br{x + z} \parder{z} {x} + \br{y + z} \parder{z} {y} = x + y$.
    
    $\dfrac {\d x} {x + z} = \dfrac {\d y} {y + z} = \dfrac {\d z} {x  + y}$.
    
    \begin{enumerate}
        \item Первую интегрируемую комбинацию можно усмотреть, если
        сложить числители и знаменатели всех трёх дробей и потом вычесть $\d x - \d y$,
        что будет пропорционально $x - y$.
        Таким образом,
        $\dfrac {\d x + \d y + \d z} {2\br{x + y + z}} = \dfrac {\d x - \d y} {x - y}$.
        $\dfrac {1} {2} \ln\vbr{x + y + z} = \ln\vbr{x - y} + A$.
        $\dfrac {x + y + z} {\br{x - y}^2} = A$.
        
        \item Вторую интегрирумую комбинацию получим следующим образом: заметим,
        что производные пропорции
        $\dvfrac {\d z - \d x}{y - z}$ и $\dfrac {\d z - \d y}{x - z}$,
        по свойству равные исходным, имеют числители и знаменатели, 
        дополняющие друг друга до полного дифференциала:
        $$\dfrac {\br{\d z - \d x} + \br{\d z - \d y}}{\br{y - z} + \br{x - z}}
        = \dfrac {-\d{\br{2z - x - y}}} {2z - x - y}
        = - \d{\ln\vbr{2z - x - y}}.$$
        Соответственно, полученное выражение по свойству производных пропорций равно
        всем ранее полученным, и в частности
        $- \d{\ln\vbr{2z - x - y}} = \dfrac {\d x - \d y} {x - y} = \d{\ln\vbr{x - y}}$,
        то есть
        $\br{2z - x - y}\br{x - y} = B$.
    \end{enumerate}
    
    Получаем общее решение
    $F\br{\dvfrac{x + y + z} {\br{x - y}^2}; \br{2z - x - y}\br{x - y}} = 0$.
\end{document}
